% =====================================================================
\chapter{Sucesiones y series numéricas}
\label{tema:5}
% =====================================================================

Una serie es una suma de infinitos términos. Como no se pueden sumar infinitos números uno tras otro, la suma se define como el límite de las sumas parciales, y la primera pregunta es si ese límite existe. Este tema repasa los límites de sucesiones, estudia la serie geométrica y la serie armónica, que son los dos modelos de comparación fundamentales, y desarrolla los criterios de convergencia: comparación, cociente, raíz, integral y Leibniz. Se siguen los capítulos~22 y~23 de Spivak~\cite{spivak}, el capítulo~10 de Apostol~\cite{apostol}, las §5.1 a~5.6 de~\cite{openstax2}, las §10.1 a~10.3 de Strang~\cite{strang} y las sesiones~94 y~95 del curso 18.01SC~\cite{mit1801}.

\section{Sucesiones}

\begin{definicion}[Límite de una sucesión]\label{def:limite-sucesion}
Una sucesión $(a_n)_{n\ge1}$ de números reales \emph{converge} a $\ell\in\R$, y se escribe $a_n\to\ell$ o $\lim_{n\to\infty}a_n=\ell$, si para cada $\varepsilon>0$ existe $N\in\N$ tal que $\abs{a_n-\ell}<\varepsilon$ para todo $n\ge N$. Se escribe $a_n\to\infty$ si para cada $M$ existe $N$ con $a_n>M$ para todo $n\ge N$. Una sucesión que no converge \emph{diverge}.
\end{definicion}

El límite, si existe, es único, y los límites de sucesiones cumplen las mismas reglas algebraicas que los de funciones (módulo~0.5): el límite de una suma, de un producto o de un cociente con denominador de límite no nulo es la suma, el producto o el cociente de los límites. También vale el criterio del encaje: si $a_n\le b_n\le c_n$ y $a_n,c_n\to\ell$, entonces $b_n\to\ell$. Si $a_n=f(n)$ y $f(x)\to\ell$ cuando $x\to\infty$, entonces $a_n\to\ell$, lo que permite usar la regla de L'Hôpital.

\begin{teorema}[Convergencia monótona]\label{teo:monotona-acotada}
Toda sucesión creciente y acotada superiormente converge a su supremo. Toda sucesión decreciente y acotada inferiormente converge a su ínfimo.
\end{teorema}

\addtocounter{definicion}{1}%

\section{Series numéricas}

\begin{definicion}[Serie]\label{def:serie}
Dada una sucesión $(a_n)_{n\ge1}$, su \emph{suma parcial} $n$-ésima es $S_n=a_1+\dots+a_n$. La serie $\sum_{n=1}^{\infty}a_n$ \emph{converge} si $(S_n)$ converge, y su \emph{suma} es $\sum_{n=1}^{\infty}a_n=\lim_{n\to\infty}S_n$. En otro caso la serie \emph{diverge}.
\end{definicion}

Si $\sum a_n$ y $\sum b_n$ convergen, también $\sum(\alpha a_n+\beta b_n)$, con suma $\alpha\sum a_n+\beta\sum b_n$, porque las sumas parciales se combinan del mismo modo. Suprimir o añadir un número finito de términos no cambia el carácter de la serie, aunque sí su suma.

\begin{teorema}[Serie geométrica]\label{teo:geometrica}
Sea $r\in\R$. La serie $\sum_{n=0}^{\infty}r^{n}$ converge si y solo si $\abs r<1$, y entonces
\[
\sum_{n=0}^{\infty}r^{n}=\frac{1}{1-r}.
\]
\end{teorema}

\begin{proposicion}[Condición necesaria]\label{prop:condicion-necesaria}
Si $\sum a_n$ converge, entonces $a_n\to0$. Equivalentemente, si $a_n\not\to0$, la serie diverge.
\end{proposicion}

La condición $a_n\to0$ es necesaria pero no suficiente, como muestra el resultado siguiente.

\begin{teorema}[La serie armónica diverge]\label{teo:armonica}
La serie $\sum_{n=1}^{\infty}\frac1n$ diverge; de hecho, $H_{2^{k}}\ge1+\frac k2$ para todo $k\in\N_0$.
\end{teorema}

\addtocounter{definicion}{1}%

\section{Series de términos no negativos}

Si $a_n\ge0$ para todo $n$, las sumas parciales forman una sucesión creciente y, por el \cref{teo:monotona-acotada}, la serie converge si y solo si las sumas parciales están acotadas. De ahí salen todos los criterios de esta sección.

\begin{teorema}[Criterio de comparación]\label{teo:comparacion-series}
Sean $0\le a_n\le b_n$ para todo $n$. Si $\sum b_n$ converge, $\sum a_n$ converge y $\sum a_n\le\sum b_n$. Si $\sum a_n$ diverge, $\sum b_n$ diverge.
\end{teorema}

\begin{teorema}[Comparación en el límite]\label{teo:comparacion-limite-series}
Sean $a_n,b_n>0$ con $\frac{a_n}{b_n}\to L\in(0,\infty)$. Entonces $\sum a_n$ y $\sum b_n$ tienen el mismo carácter.
\end{teorema}

\begin{teorema}[Criterio de la integral]\label{teo:criterio-integral}
Sea $f\colon[1,\infty)\to\R$ positiva y decreciente y sea $a_k=f(k)$. Para todo $n\ge1$,
\[
\int_1^{n+1}f(x)\dif x\le\sum_{k=1}^{n}a_k\le a_1+\int_1^{n}f(x)\dif x .
\]
En consecuencia, $\sum a_k$ converge si y solo si $\int_1^{\infty}f$ converge.
\end{teorema}

\addtocounter{figure}{1}%

\begin{corolario}[Series $p$]\label{cor:series-p}
La serie $\sum_{n=1}^{\infty}\frac{1}{n^{p}}$ converge si y solo si $p>1$. Además, $\ln(n+1)\le H_n\le1+\ln n$ para todo $n\in\N$.
\end{corolario}

Así, la serie armónica diverge, pero muy despacio: $H_{1000}$ está entre $\ln1001\approx6{,}91$ y $1+\ln1000\approx7{,}91$. En cambio, $\sum\frac{1}{n^{2}}$ converge, y su suma es menor que $1+\int_1^\infty x^{-2}\dif x=2$ (su valor exacto, $\frac{\pi^{2}}{6}\approx1{,}645$, se obtiene con otras técnicas).

\section{Criterios del cociente y de la raíz}

\begin{teorema}[Criterio del cociente]\label{teo:cociente}
Sea $a_n>0$ con $\frac{a_{n+1}}{a_n}\to L$. Si $L<1$, $\sum a_n$ converge; si $L>1$ (incluido $L=\infty$), diverge. Si $L=1$, el criterio no decide.
\end{teorema}

\begin{teorema}[Criterio de la raíz]\label{teo:raiz}
Sea $a_n\ge0$ con $\sqrt[n]{a_n}\to L$. Si $L<1$, $\sum a_n$ converge; si $L>1$, diverge. Si $L=1$, el criterio no decide.
\end{teorema}

La demostración del criterio del cociente es análoga y se pide en el \cref{eje:5-cociente-dem}. Ambos criterios comparan con una serie geométrica, y por eso no deciden en el caso $L=1$: tanto $\sum\frac1n$, divergente, como $\sum\frac{1}{n^{2}}$, convergente, tienen $L=1$. Son especialmente útiles con factoriales y potencias $n$-ésimas.

\addtocounter{definicion}{1}%

\section{Series alternadas y convergencia absoluta}

\begin{definicion}[Convergencia absoluta y condicional]\label{def:absoluta}
La serie $\sum a_n$ converge \emph{absolutamente} si $\sum\abs{a_n}$ converge. Si $\sum a_n$ converge pero $\sum\abs{a_n}$ diverge, converge \emph{condicionalmente}.
\end{definicion}

\begin{teorema}\label{teo:absoluta-series}
Toda serie absolutamente convergente es convergente, y $\bigl|\sum a_n\bigr|\le\sum\abs{a_n}$.
\end{teorema}

\begin{teorema}[Criterio de Leibniz]\label{teo:leibniz}
Sea $(b_n)$ decreciente con $b_n\to0$ (luego $b_n\ge0$). La serie alternada $\sum_{n=1}^{\infty}(-1)^{n+1}b_n=b_1-b_2+b_3-\cdots$ converge, y su suma $S$ cumple $\abs{S-S_n}\le b_{n+1}$ para todo $n$.
\end{teorema}

\addtocounter{definicion}{1}%

\begin{observacion}[Reordenaciones]
Una serie absolutamente convergente puede sumarse en cualquier orden con el mismo resultado. Con una serie condicionalmente convergente no ocurre así: un teorema de Riemann afirma que sus términos pueden reordenarse para obtener cualquier suma prefijada~\cite[cap.~23]{spivak}. El ejercicio~25 de la §10.3 de Strang~\cite{strang} muestra una reordenación de la serie armónica alternada cuya suma es $\frac12\ln2$.
\end{observacion}

\addtocounter{definicion}{1}%

\addtocounter{definicion}{1}%

